% ===================================================================
% FST-IV: Cosmological Stability -- a Collector Paper
%   for CRM Series, Saturation Theorem, QG-CRM and Dark Energy
% Target journal: Foundations of Physics / Universe / JCAP
% Status: v0.1 SKELETON (programmatic survey of the cosmological scale
%                       under the FST-UCU architecture)
% Author: Lukas Geiger, Bernau, Germany
% ===================================================================

\documentclass[11pt,a4paper]{article}

\usepackage[utf8]{inputenc}
\usepackage[T1]{fontenc}
\usepackage{amsmath,amssymb,amsthm}
\usepackage{graphicx}
\usepackage{hyperref}
\usepackage{booktabs}
\usepackage[margin=2.4cm]{geometry}
\usepackage{enumitem}
\usepackage{xcolor}

\newtheorem{axiom}{Axiom}
\newtheorem{definition}{Definition}
\newtheorem{lemma}{Lemma}
\newtheorem{proposition}{Proposition}
\newtheorem{conjecture}{Conjecture}
\theoremstyle{remark}
\newtheorem*{remark}{Remark}

\title{Functional Stability Theory IV:\\
Cosmological Stability\\
\large\textit{A Collector Paper for CRM, Saturation,
QG-CRM and Dark Energy under the FST-UCU Architecture}}

\author{Lukas Geiger\\
\textit{Independent Researcher, Bernau, Germany}\\
ORCID: \url{https://orcid.org/0009-0005-7296-1534}}

\date{May 2026}

\begin{document}
\maketitle

\begin{abstract}
\noindent
This paper is the \emph{cosmological collector slot} (FST-IV) of the
Functional Stability Theory (FST) programme, sibling to FST-I
(thermodynamic stability of fundamental parameters), FST-II (chemical
stability and autocatalytic selection), and FST-III (biological
stability and Nash frustration). The cosmological scale of the FST
programme is already addressed by several published companion papers:
the Cooperative Renormalization Model (CRM) Papers I-IV, the Saturation
Theorem (CRM-V), the inflationary QG-CRM line (CRM-VI), and the
Dark-Energy variational bound (FST-DE). This collector does \emph{not}
replace any of these. It (i) embeds each component in the FST-UCU
architecture (universal game, MEPP axiom, three-lemma endgame UCU1-3,
RG bridge), (ii) clarifies which UCU step is supplied by which paper,
(iii) records the open sub-problems honestly (L4 RG-matching from
quadratic gravity to a Hu--Sawicki IR effective theory, the
Convexity-Screening No-Go for linear $R^2$, and the Pattern-A status
of several reductions), and (iv) includes a mathematical reformulation
of CRM as a one-parameter Boost subgroup of $\mathrm{PSL}_2(\mathbb{R})$
acting on hyperbolic geometry, opening a Flow-Zeta classification
parallel to Riemann (arithmetic), Selberg (geometric-compact), and the
non-compact dynamical zoo. \textbf{This paper is an honest survey, not
a new theorem.} The programmatic umbrella is FST -- A Programmatic
Hub~\cite{GeigerFST-Hub-2026}.
\end{abstract}

\textbf{Keywords:} cosmology, dark energy, cosmological constant,
scalar-tensor gravity, quadratic gravity, $f(R)$ gravity, inflation,
Saturation Theorem, Pattern A, functional stability theory, FST,
Flow-Zeta, hyperbolic geometry.

% ===================================================================
\section{Introduction}
\label{sec:intro}
% ===================================================================

The FST programme posits that emergent stability across scales obeys
a single conditional architecture: a positive functional, a strict
convexity property (UCU1), an existence/coercivity property (UCU2),
and a quantitative quadratic lower bound (UCU3) on a closed convex
feasibility set $K \subset X$, together with a Renormalization-Group
bridge connecting adjacent scales. The application companions FST-I,
FST-II and FST-III instantiate this on the particles, chemistry and
biology scales respectively. The cosmological scale -- which this
collector addresses -- already hosts a richer set of pre-existing
papers, and our task is not to reinvent them but to \emph{place} them
in the same UCU/RG language and audit their honest status.

% ===================================================================
\section{The Cosmological Scale at a Glance}
\label{sec:components}
% ===================================================================

Table~\ref{tab:cosmo-components} lists the cosmological FST components.

\begin{table}[h]
\centering
\renewcommand{\arraystretch}{1.2}
\small
\begin{tabular}{p{3.0cm}p{3.0cm}p{5.0cm}p{2.5cm}}
\toprule
\textbf{Component} & \textbf{FST-UCU role} & \textbf{Content} &
\textbf{Concept-DOI} \\
\midrule
CRM Papers I-IV   & game-theoretic foundation; UCU1+2 candidate     & spieltheoretic CRM, MOND unification, scalaron Lagrangian $R+\gamma R^2$, vector sector & 10.5281/zenodo.18728936 \\
CRM-V (Saturation)& \emph{Universal} tanh normal form               & axioms A-D, Abel equation, tanh as unique IR attractor  & 10.5281/zenodo.19243234 \\
CRM-VI (QG-CRM)   & UV completion; L4 RG-matching open              & quadratic gravity inflation, Hu--Sawicki IR bridge      & 10.5281/zenodo.19352449 (draft) \\
FST-DE            & Variational bound; Convexity-Screening No-Go    & $\Phi[\rho]$ minimum yields $\Lambda_\text{eff}\sim H_0^2/\ln(M_\text{Pl}/H_0)$; gravitational slip $\eta=5/4$ & 10.5281/zenodo.19036235 \\
\bottomrule
\end{tabular}
\caption{Cosmological FST components, their UCU role and persistent
identifiers. The collector embeds each in the FST architecture; it
does not replace any of them.}
\label{tab:cosmo-components}
\end{table}

% ===================================================================
\section{The Universal Game at Cosmological Scale}
\label{sec:game}
% ===================================================================

\textbf{TODO v0.2 -- v1.0:} Identify $(x_L, G_L, D_L)$ for the
cosmological scale: $x_L$ is the curvature/density configuration,
$G_L$ is the spieltheoretic game from CRM Paper~I (player = vacuum
field vs.\ blase), payoff = integrated entropy production over the
expansion history. MEPP axiom: cosmological evolution maximises
entropy production along open trajectories. RG bridge: connects the
$R + \gamma R^2$ scalaron sector to the IR effective $f(R)$.

% ===================================================================
\section{UCU1/UCU2/UCU3 Status across the Cosmological Components}
\label{sec:ucu-status}
% ===================================================================

\textbf{TODO v0.2 -- v1.0:} Detailed UCU status table.
Schematic:

\begin{itemize}
\item \textbf{UCU1 (strict convexity):} verified for $\Phi[\rho]$
  (FST-DE), conditional for the CRM tanh-saturation potential.
\item \textbf{UCU2 (existence):} Tonelli direct method on
  reflexive Banach spaces; available throughout the cosmological
  sector.
\item \textbf{UCU3 local:} achieved for $\Phi[\rho]$ at the de~Sitter
  fixed point.
\item \textbf{UCU3 global (L4 RG-matching):} \emph{open}. This is
  the structural gap CRM-VI documents: lifting the variational bound
  from $\Phi[\rho]$ through a four-step RG match to an effective
  $f(R)$ theory at the IR fixed point.
\end{itemize}

% ===================================================================
\section{Mathematical Reformulation: CRM as Flow-Zeta on Hyperbolic
Geometry}
\label{sec:flow-zeta}
% ===================================================================

\textbf{TODO v0.2 -- v1.0:} Integrate the material from the working
note \texttt{CRM\_AS\_FLOW\_ZETA.md}. Key structural identifications:

\begin{itemize}
\item The CRM composition law $x \oplus y = (x+y)/(1+xy)$ on
  $(-1,1)$ is isomorphic to $(\mathbb{R}, +)$ via $\tanh^{-1}$ and
  realises a one-parameter Boost subgroup of $\mathrm{PSL}_2(\mathbb{R})$.
\item This places CRM in the three-class Zoo taxonomy alongside
  Riemann (arithmetic), Selberg (geometric-compact) and the
  non-compact dynamical zeta family (Ruelle-Anosov, BH quasinormal
  modes).
\item The HP-YES claim for CRM is \emph{conditional} on the SGE
  conjecture (Zoo Master, CORE/zetazoo); it is a classification
  statement, not an independent proof.
\end{itemize}

% ===================================================================
\section{Pattern-A Audit}
\label{sec:pattern-a}
% ===================================================================

\textbf{TODO v0.2 -- v1.0:} Apply the Pattern-A negative form
(positivity reformulation alone is in general only a translation of
language). The cosmological components contain two clear Pattern-A
reductions:
\begin{enumerate}
\item \textbf{FST-DE Convexity-Screening No-Go:} strict convexity of
  $\Phi$ admits a unique minimiser, so linear $R+\gamma R^2$ cannot
  produce density-dependent screening. The bridge to Hu--Sawicki is
  a \emph{conditional structural reduction}, not a closure.
\item \textbf{CRM-VI L4 RG-matching:} mapping the QG-CRM UV theory
  onto a Hu--Sawicki IR $f(R)$ is a sub-conditional reduction with
  two explicitly identified open sub-conditions (B1 trajectory
  existence, B2 double-logarithm matching).
\end{enumerate}

% ===================================================================
\section{Honest Status / Open Sub-problems}
\label{sec:honest}
% ===================================================================

\begin{description}[leftmargin=2em,style=nextline]
\item[What this collector delivers.] An FST-UCU embedding of the
existing cosmological components (CRM I-IV, CRM-V, CRM-VI, FST-DE);
a status map of UCU1/UCU2/UCU3 across them; the Pattern-A audit of
the two reductions above; the Flow-Zeta classification of CRM
within the Zeta Zoo.
\item[What this collector does NOT deliver.] A closure of the L4
RG-matching (CRM-VI). A non-conditional Convexity-Screening
resolution (FST-DE). A non-SGE-conditional HP-YES proof for CRM.
A new dark-energy theorem.
\item[Falsification criterion.] If UCU1 or UCU2 fails on the
natural cosmological feasibility set, the collector and the
underlying components require structural revision. UCU3-global
failure is not falsification; it is the locator of the open
sub-problem.
\end{description}

% ===================================================================
\section*{Acknowledgements}
% ===================================================================
This paper synthesises work from the FST-DE Dark Energy line and
the CRM series (Papers I-IV, CRM-V, CRM-VI). Computations supporting
the FST-DE MCMC and the Hu--Sawicki phase-space fit were performed
on \texttt{ellmos-services} and the Mac Studio compute host.

% ===================================================================
\section*{AI Disclosure (Level L3)}
% ===================================================================
This paper was prepared with substantial AI assistance (Anthropic
Claude, OpenAI GPT, Google Gemini variants). AI tools were used for
cross-paper synthesis between FST-DE and the CRM line, structural
editing, status-map generation, Flow-Zeta classification literature
cross-referencing, and bilingual EN/DE preparation. All scientific
claims, the UCU architecture, the Pattern-A audit and the
falsification criterion are the author's own and are subject to
standard authorial responsibility. No section was written
autonomously by AI.

% ===================================================================
\begin{thebibliography}{99}

\bibitem{GeigerFST-Hub-2026}
L.~Geiger,
\emph{Functional Stability Theory --- A Programmatic Hub:
Universal Convexity-Uniqueness Across Scales},
Zenodo: \texttt{10.5281/zenodo.20130499}, 2026.

\bibitem{GeigerCRM2026}
L.~Geiger,
\emph{Cooperative Renormalization Model -- Papers I-IV},
Zenodo: \texttt{10.5281/zenodo.18728936}, 2026.

\bibitem{GeigerCRMV2026}
L.~Geiger,
\emph{The Saturation Theorem (CRM Paper V)},
Zenodo: \texttt{10.5281/zenodo.19243234}, 2026.

\bibitem{GeigerCRMVI2026}
L.~Geiger,
\emph{Quantum Quadratic Gravity Completion of CRM (CRM-VI, draft)},
Zenodo: \texttt{10.5281/zenodo.19352449}, 2026.

\bibitem{GeigerFST-DE-2026}
L.~Geiger,
\emph{Dark Energy as Residual Vacuum Free Energy: A Thermodynamic
Bound on the Cosmological Constant},
Zenodo: \texttt{10.5281/zenodo.19036235}, 2026.

\bibitem{LiuQuintinAfshordi2026}
J.~Liu, J.-L.~Quintin, N.~Afshordi,
\emph{Ultraviolet Completion of the Big Bang in Quadratic Gravity},
Phys.\ Rev.\ Lett.\ \textbf{136}, 111501 (2026).

\bibitem{HuSawicki2007}
W.~Hu, I.~Sawicki,
\emph{Models of $f(R)$ cosmic acceleration that evade solar-system
tests}, Phys.\ Rev.\ D \textbf{76}, 064004 (2007).

\bibitem{Stelle1977}
K.~S.~Stelle,
\emph{Renormalization of higher-derivative quantum gravity},
Phys.\ Rev.\ D \textbf{16}, 953 (1977).

\bibitem{Starobinsky1980}
A.~A.~Starobinsky,
\emph{A new type of isotropic cosmological models without
singularity}, Phys.\ Lett.\ B \textbf{91}, 99 (1980).

\bibitem{DyatlovZworski2019}
S.~Dyatlov, M.~Zworski,
\emph{Mathematical Theory of Scattering Resonances},
Amer.\ Math.\ Soc., 2019.

\end{thebibliography}

\end{document}
